\documentclass{exam}
\usepackage{amsmath, amssymb}

\pagestyle{headandfoot}
\firstpageheadrule
\runningheadrule
\firstpageheader{Prof. Tahvildar-Zadeh \\ Linear Algebra}{Homework 5}{Aadhithya Saravanan}
\runningheader{Linear Algebra \\ Homework 5}{}{Saravanan}
\firstpagefooter{}{}{}
\runningfooter{}{\thepage}{}

\printanswers

\begin{document}

\underline{Exercise 2.3.3}
\newline
\begin{parts}
    \part To compute $[U]_\beta^\gamma$, we need to take the standard basis vectors of $P_2(\mathbb{R})$, apply $U$ to them, and write them using the standard basis vectors of $\mathbb{R}^3$. The vectors are $\{1,x,x^2\}$. $U(1)=(1,0,1)=e_1+e_3$, $U(x)=(1,0,-1)=e_1-e_3$, and $U(x^2)=(0,1,0)=e_2$. Arranging these in matrix form, we have
    $\begin{bmatrix}
        1 & 1 & 0 \\
        0 & 0 & 1 \\
        1 & -1 & 0
    \end{bmatrix}$, which is $[U]_\beta^\gamma$. Doing the same for $[T]_\beta^\beta$, we apply $T$ to each vector in $\{1,x,x^2\}$. $T(1)=2=2e_1$, $T(x)=3+3x=3e_1+3e_2$, and $T(x^2)=6x+4x^2=6e_2+4e_3$. Arranging this in matrix form, we have
    $\begin{bmatrix}
        2 & 3 & 0 \\
        0 & 3 & 6 \\
        0 & 0 & 4
    \end{bmatrix}$, which is $[T]_\beta^\beta$. To do the same for $[UT]_\beta^\gamma$, we apply $T$ and then $U$ to each vector in $\{1,x,x^2\}$. $U(T(1))=(2,0,2)=2e_1+2e_3$, $U(T(x))=(6,0,0)=6e_1$, and $U(T(x^2))=6x+4x^2=(6,4,-6)=6e_1+4e_2-6e_3$. Arranging this in matrix form, we have
    $\begin{bmatrix}
        2 & 6 & 6 \\
        0 & 0 & 4 \\
        2 & 0 & -6
    \end{bmatrix}$, which is $[UT]_\beta^\gamma$. Theorem 2.11 says that $[UT]_\beta^\gamma=[U]_\beta^\gamma[T]_\beta^\beta$, which we can verify by matrix multiplication. Since
    $\begin{bmatrix}
        1 & 1 & 0 \\
        0 & 0 & 1 \\
        1 & -1 & 0
    \end{bmatrix}
    \begin{bmatrix}
        2 & 3 & 0 \\
        0 & 3 & 6 \\
        0 & 0 & 4
    \end{bmatrix}=
    \begin{bmatrix}
        2 & 6 & 6 \\
        0 & 0 & 4 \\
        2 & 0 & -6
    \end{bmatrix}$, we computed $[UT]_\beta^\gamma$ correctly.
    \part $[h(x)]_\beta$ can be computed by writing $h(x)$ in the standard basis vectors of $P_2(\mathbb{R})$, which are $\{1,x,x^2\}$. So, $[h(x)]_\beta=3e_1-2e_2+e_3=
    \begin{bmatrix}
    3 \\
    -2 \\
    1
    \end{bmatrix}$. To compute $[U(h(x))]_\gamma$, we need to apply $U$ to $h(x)$ and then write the result in the standard basis vectors of $\mathbb{R}^3$. $U(h(x))=(1,1,5)=e_1+e_2+5e_3$. So, $[U(h(x))]_\gamma=
    \begin{bmatrix}
        1 \\
        1 \\
        5
    \end{bmatrix}$. Theorem 2.14 says that $[U(h(x))]=[U]_\beta^\gamma[h(x)]_\beta$. Since $
    \begin{bmatrix}
        1 & 1 & 0 \\
        0 & 0 & 1 \\
        1 & -1 & 0
    \end{bmatrix}
    \begin{bmatrix}
        3 \\
        -2 \\
        1
    \end{bmatrix}=
    \begin{bmatrix}
        1 \\
        1 \\
        5
    \end{bmatrix}$, we computed $[U(h(x))]$ correctly.
    \newline
\end{parts}

\underline{Exercise 2.3.4.d}
\newline
\newline
Exercise 2.2.5.d defines $T:P_2(\mathbb{R})\to\mathbb{R}$ as $T(f(x))=f(2)$. Applying $T$ to the basis vectors of $P_2(\mathbb{R})$, which are $\{1,x,x^2\}$, gives $T(1)=1$, $T(x)=2$, and $T(x^2)=4$. So, letting $\beta$ be the standard basis of $P_2(\mathbb{R})$ and $\gamma$ be the standard basis of $\mathbb{R}$, we have $[T]_\beta^\gamma=
\begin{bmatrix}
    1 & 2 & 4
\end{bmatrix}$. Now, to compute $[T(f(x))]_\gamma$ when $f(x)=6-x+2x^2$, we can use Theorem 2.14, which states that $[T(f(x))]_\gamma=[T]_\beta^\gamma[f(x)]_\beta$. Writing $f(x)$ in the standard basis vectors of $P_2(\mathbb{R})$, we have $[f(x)]_\beta=6e_1-e_2+2e_3=
\begin{bmatrix}
    6 \\
    -1 \\
    2
\end{bmatrix}$. Then, we can compute $[T(f(x))]_\gamma=[T]_\beta^\gamma[f(x)]_\beta=
\begin{bmatrix}
    1 & 2 & 4
\end{bmatrix}
\begin{bmatrix}
    6 \\
    -1 \\
    2
\end{bmatrix}
=12$.
\newline

\underline{Exercise 2.4.2}
\newline
\begin{parts}
    \part $T$ is not invertible.  By the corollary to Theorem 2.18, $T$ is invertible if and only if $[T]_\beta^\gamma$ is invertible. In this case $\beta$ is the standard basis of $\mathbb{R}^2$ and $\gamma$ is the standard basis of $\mathbb{R}^3$. However, these bases have different dimensions, so $[T]_\beta^\gamma$ is not square. Then, $[T]_\beta^\gamma$ is also not invertible, so $T$ is not invertible.
    \part $T$ is not invertible.  By the corollary to Theorem 2.18, $T$ is invertible if and only if $[T]_\beta^\gamma$ is invertible. In this case $\beta$ is the standard basis of $\mathbb{R}^2$ and $\gamma$ is the standard basis of $\mathbb{R}^3$. However, these bases have different dimensions, so $[T]_\beta^\gamma$ is not square. Then, $[T]_\beta^\gamma$ is also not invertible, so $T$ is not invertible.
    \part $T$ is invertible. By the corollary to Theorem 2.18, $T$ is invertible if and only if $[T]_\beta^\gamma$ is invertible. $[T]_\beta^\gamma=
    \begin{bmatrix}
        3 & 0 & -2 \\
        0 & 1 & 0 \\
        3 & 4 & 0
    \end{bmatrix}$. The determinant of the matrix is 6, which is nonzero, so the matrix is invertible. Therefore, $T$ is invertible.
    \part $T$ is not invertible.  By the corollary to Theorem 2.18, $T$ is invertible if and only if $[T]_\beta^\gamma$ is invertible. In this case $\beta$ is the standard basis of $P_3(\mathbb{R})$ and $\gamma$ is the standard basis of $P_2(\mathbb{R})$. However, these bases have different dimensions, so $[T]_\beta^\gamma$ is not square. Then, $[T]_\beta^\gamma$ is also not invertible, so $T$ is not invertible.
    \newline
\end{parts}

\underline{Exercise 2.4.3.a}
\newline
\newline
By the corollary of Theorem 2.19, $F^3$ is isomorphic to $P_3(F)$ if and only if $\dim(P_3(F))=3$. $\{1,x,x^2,x^3\}$ is a basis for $P_3(F)$, so $\dim(P_3(F))=4$. Therefore, $F^3$ is not isomorphic to $P_3(F)$.
\newline

\underline{Exercise 2.4.3.c}
\newline
\newline
By Theorem 2.19, $M_{2\times2}(\mathbb{R})$ is isomorphic to $P_3(\mathbb{R})$ if and only if $\dim(M_{2\times2}(\mathbb{R}))=\dim(P_3(\mathbb{R}))$. A basis for $M_{2\times2}(\mathbb{R})$ is $\{
    \begin{bmatrix}
        1 & 0 \\
        0 & 0
    \end{bmatrix}
    ,
    \begin{bmatrix}
        0 & 1 \\
        0 & 0
    \end{bmatrix}
    ,
    \begin{bmatrix}
        0 & 0 \\
        1 & 0
    \end{bmatrix}
    ,
    \begin{bmatrix}
        0 & 0 \\
        0 & 1
    \end{bmatrix}
\}$, so $\dim(M_{2\times2}(\mathbb{R}))=4$. A basis for $P_3(\mathbb{R})$ is $\{1,x,x^2,x^3\}$, so $\dim(P_3(\mathbb{R}))=4$. Therefore, $\dim(M_{2\times2}(\mathbb{R}))=\dim(P_3(\mathbb{R}))$ and $M_{2\times2}(\mathbb{R})$ is isomorphic to $P_3(\mathbb{R})$.
\newline

\underline{Exercise 2.4.7.a}
\newline
\newline
Assume $A$ is invertible. If $A^2=0$, we know $AA=0$. We can right multiply by $A^{-1}$ to get $AAA^{-1}=0A^{-1}$. Then, we have $AI=0$, which implies $A=0$. However, the determinant of the zero matrix is 0, and $A$ is not invertible. Therefore, $A$ is not invertible.

\end{document}